\section*{ВВЕДЕНИЕ}
\addcontentsline{toc}{section}{ВВЕДЕНИЕ}

Формальные модели вычислений, к которым относится продукционная система Маркова, занимают важное место в теории алгоритмов и математической логике. Такие системы служат базисом для понимания фундаментальных принципов обработки информации, построения формальных языков и доказательства вычислимости различных функций. Разработанная Андреем Андреевичем Марковым в середине XX века нормальная алгоритмическая система (или система Маркова) стала одной из первых формальных моделей вычислений, эквивалентных машине Тьюринга по вычислительной мощности.

Продукционная система Маркова представляет собой совокупность правил переписывания, последовательно применяемых к строке символов. Алгоритм функционирует на основе строгой последовательности замены подстрок в тексте, пока не будет достигнуто завершающее состояние. Несмотря на внешнюю простоту, такие системы позволяют описывать и моделировать широкий класс алгоритмов, а также служат хорошим инструментом для изучения понятий вычислимости и формальных грамматик.

Практическое значение моделирования системы Маркова заключается в возможности демонстрации процессов работы формальных алгоритмов, обучения студентов принципам функционирования систем подстановок, а также проверки корректности и полноты наборов правил. Программная реализация симулятора системы Маркова позволяет не только пошагово проследить выполнение алгоритма, но и собрать статистику его работы, включая количество применённых правил и изменение длины строки на каждом этапе.

Современные средства разработки предоставляют широкие возможности для визуализации работы подобных систем. Использование языка программирования Python в сочетании с библиотеками \texttt{tkinter} и \texttt{customtkinter} позволяет создать наглядный и удобный графический интерфейс, обеспечивающий простоту взаимодействия с пользователем. В разработанной программе реализованы функции пошагового выполнения, автоматического запуска алгоритма, сохранения логов преобразований и вывода статистики выполнения.

\emph{Цель настоящей работы} — разработка программной реализации симулятора продукционной системы Маркова с графическим интерфейсом, обеспечивающего пошаговое и автоматическое выполнение алгоритма, логирование всех этапов и наглядное представление процесса преобразований.

Для достижения поставленной цели необходимо решить следующие \emph{задачи}:
\begin{itemize}
	\item провести анализ теоретических основ продукционных систем Маркова;
	\item разработать архитектуру программы для моделирования системы Маркова;
	\item реализовать загрузку правил из текстового файла и ручной ввод правил через интерфейс;
	\item обеспечить выполнение алгоритма как в пошаговом режиме (с анимацией), так и в автоматическом режиме;
	\item реализовать сохранение журнала преобразований и формирование статистики выполнения;
	\item протестировать работу симулятора и продемонстрировать примеры работы системы.
\end{itemize}

\emph{Структура и объём работы.} Отчёт состоит из введения, четырёх разделов основной части, заключения, списка использованных источников и приложений. Текст курсовой работы равен \formbytotal{lastpage}{страниц}{е}{ам}{ам}.

\emph{Во введении} сформулирована цель и задачи исследования, дана общая характеристика темы и описана структура работы.

\emph{В первом разделе} приводится анализ предметной области, рассматриваются теоретические основы систем Маркова, их структура и принципы функционирования.

\emph{Во втором разделе} описано техническое задание на разработку программного симулятора, сформулированы требования к функциональности и интерфейсу программы.

\emph{В третьем разделе} рассматривается процесс технического проектирования: архитектура программы, взаимодействие модулей, структура классов и организация данных.

\emph{В четвёртом разделе} представлены результаты реализации симулятора, проведено тестирование и приведены примеры работы программы.

В заключении подведены итоги работы, сделаны выводы о соответствии поставленной цели и достигнутых результатах.

В приложении~А приведены иллюстрации и схемы архитектуры программы.  
В приложении~Б представлены фрагменты исходного кода разработанного приложения.